\documentclass[11pt,reqno]{amsart}
\usepackage[T1]{fontenc}
\usepackage[utf8]{inputenc}
\usepackage{lmodern,microtype,mathtools,amssymb,booktabs,array}
\usepackage[margin=1in]{geometry}
\usepackage[hidelinks]{hyperref}
\usepackage{enumitem}
\numberwithin{equation}{section}
\newtheorem{theorem}{Theorem}[section]
\newtheorem{lemma}[theorem]{Lemma}
\newtheorem{proposition}[theorem]{Proposition}
\newtheorem{corollary}[theorem]{Corollary}
\theoremstyle{definition}
\newtheorem{definition}[theorem]{Definition}
\theoremstyle{remark}
\newtheorem{remark}[theorem]{Remark}
\newtheorem*{theoremA}{Theorem A}
\newtheorem*{theoremB}{Theorem B}
\newtheorem*{theoremC}{Theorem C}
\newcommand{\R}{\mathbb R}
\newcommand{\C}{\mathbb C}
\newcommand{\Q}{\mathbb Q}
\newcommand{\HH}{\mathbb H}
\newcommand{\F}{\mathbb F}
\newcommand{\Sym}{\mathbb S}
\newcommand{\Herm}{\mathcal H}
\newcommand{\rank}{\operatorname{rank}}
\newcommand{\tr}{\operatorname{tr}}
\newcommand{\dist}{\operatorname{dist}}
\newcommand{\supp}{\operatorname{supp}}
\newcommand{\spanof}{\operatorname{span}}
\newcommand{\ri}{\operatorname{ri}}
\newcommand{\Id}{\mathrm{Id}}
\newcommand{\ip}[2]{\langle #1,#2\rangle}
\newcommand{\norm}[1]{\lVert #1\rVert}
\newcommand{\one}{\mathbf 1}
\DeclareMathOperator{\den}{den}
\title[Rank universality and field-dependent tangency]{Rank universality and field-dependent tangency of symmetric cones}
\date{November 12, 2026\quad\textnormal{--- consolidated pre-submission research draft}}
\subjclass[2020]{Primary 52A20, 90C22; Secondary 14P10, 17C20, 15B48}
\keywords{Symmetric cone, sharp error bound, facial reduction, singularity degree, central path, positive semidefinite matrix, valuation}
\begin{document}
\begin{abstract}
We classify the optimal H\"older error exponents of affine sections of irreducible symmetric cones jointly with minimum facial-reduction depth at fixed cone rank and feasible rank. The spectrum is independent of the irreducible cone type and has a finite arithmetic description by harmonic midpoint weights and two peeling recurrences. A one-switch construction realizes every permitted depth without increasing rank. Over complex and quaternionic Hermitian positive semidefinite cones, a survivor-table refinement classifies the entire central valuation spectrum, its multiplicities, the relative position of central and facial flags, and the greedy facial rank-drop profile. This finer geometry remains field-dependent even when the central spectrum is not prescribed. For every $m\ge1$, the exponent $5/27$ and profile $(m,1,2,2)$ at nullity $m+5$ are attainable over $\F\in\{\R,\C,\HH\}$ at positive feasible rank $r$ exactly when $r\dim_{\R}\F\ge2$. The exclusion of all real rank-one realizations follows from a harmonic denominator restriction and an off-face positive-circuit obstruction. Explicit rational pencils prove all converse assertions. The proofs separate the numerical order of tangency from the finer geometry of facial exposure.
\end{abstract}
\maketitle
\tableofcontents
\section{Introduction}
An affine space may meet the positive semidefinite cone in a single matrix without meeting it transversely. In that situation, a small error in the equations can permit a much larger displacement of a positive semidefinite matrix. The optimal power relating these quantities measures the order of tangency of the affine space to the cone. A different invariant is discrete: facial reduction uncovers the smallest face containing the section, and its minimum length records how many successive exposures are needed. This paper determines which analytic and facial invariants are compatible, and where that compatibility depends on the underlying scalar field.

The elementary example already distinguishes curved from polyhedral geometry. At $X_*=\operatorname{diag}(1,0)$, measurement of the two diagonal entries leaves an offdiagonal perturbation of size $\sqrt{\delta}$ possible when the second diagonal is $\delta$. The sharp recovery exponent is $1/2$, although the identity-shift maximum-determinant center is $\operatorname{diag}(1+t,t)$ and moves only linearly in $t$. Thus the displacement of the center is not the sharp error modulus. Conversely, the eigenvalue orders of that center contain geometric information that is absent from a single worst-case bound.

Facial reduction was developed by Borwein and Wolkowicz \cite{BW81}. Sturm's semidefinite error bound \cite{Sturm00} relates a guaranteed H\"older exponent to the singularity degree. Louren\c co's theory of amenable cones and facial residual functions \cite{Lourenco21} supplies error bounds on symmetric cones without constraint qualifications. These results provide estimates valid for broad classes of sections. The question considered here is also a realization question: at a fixed cone rank and feasible rank, which \emph{sharp} exponents occur at a specified minimum facial depth, and which complete sequences of rank drops can accompany them?

Our first result gives a finite arithmetic answer independent of the irreducible symmetric-cone type. The relevant size is Jordan rank, not the dimension of the ambient real vector space. The allowed central orders form a strictly increasing list
\[
0=w_0<w_1<\cdots<w_k=1
\]
whose internal entries satisfy midpoint relations $2w_i=w_a+w_b$ with $a<i<b$. Two integer recurrences on this list give a lower and an upper facial depth. Every integer between them is attained at the original rank, with any admissible leading half-weight and any positive multiplicities of the central grades. The realization follows the upper recurrence for an initial segment and then switches once to the lower recurrence. Selecting this family of schedules, rather than trying to realize every schedule over the real field, is essential.

Our second result retains the complete facial information over $\C$ and $\HH$. After a congruence preserving central orders, the relative position of the central flag and the greedy facial flag is an integer table. Two local inequalities on that table are necessary and sufficient. The construction uses the real and imaginary components of a Hermitian entry independently: one carries a tight midpoint relation and the other releases a facial equation after an endpoint has disappeared. A self-release rule makes this possible for arbitrary, not necessarily prefix, flags and for grades split across several steps. No additional null direction is introduced.

The real and complex classifications cannot be identified by forgetting the central spectrum. Our third result exhibits an infinite family for which the same exponent and complete rank-drop profile are attainable over $\C$ but not over $\R$ at source rank one. Specifically, exponent $5/27$, nullity $m+5$, and profile $(m,1,2,2)$ are attainable precisely when the source-to-null entry space has at least two real coordinates. The exclusion permits arbitrary dense kernels, arbitrary real coefficients, arbitrary rotations, and every replacement central spectrum. A denominator restriction first reduces all possible central orders to two alternatives. One is excluded by the table inequalities; the other exhausts a one-dimensional real off-face space before the next positive grade can be exposed. A six-generator complex pencil attains the missing data. A second positive real source direction gives a real realization.

These results distinguish three levels of geometry. Sharp exponent and minimum depth are rank-universal across irreducible cones. Central spectral orders can also be prescribed universally when only the length of the facial chain is fixed. Complete facial profiles, however, can detect the scalar field and positive source rank even after the central orders are erased. The separation is an explicit boundary to universality, not a contradiction of its coarser form.

\subsection{Relation to central-path and tropical methods}
The regularization used throughout is the external maximum-determinant feasibility path
\[
C_t=\arg\max\{\log\det X:X\succ0,\ X-X_*-tI\in\ker A\}.
\]
Sremac, Woerdeman, and Wolkowicz \cite{SWW21}, especially their results on eigenvalue decay and singularity degree, are the closest central-path predecessor. We do not introduce this path. We identify a fixed graded order chart for it, recover the sharp full-ambient modulus from its restricted barrier Hessian, and classify the simultaneous realizable orders and facial data. The general algebraicity and Puiseux analysis of SDP central paths is likewise established; see Basu and Mohammad-Nezhad \cite{BMN22}. Their standard optimization central paths should not be conflated with the specified feasibility path here. Effective semialgebraic error estimates \cite{BMN24} provide another important background, but rationality and effective decidability alone do not imply the realization statements proved below.

The valuation inequality for a positive matrix, $w_a+w_b\le2v_{ab}$, has a direct tropical antecedent in Yu's description of the tropical positive semidefinite cone \cite{Yu15}. Generic tropical spectrahedra and stochastic games are related in \cite{AGS18}; signed tropical polarity is studied in \cite{AAGS25}. Midpoint networks also occur in mediated sets and geometric-mean conic representations \cite{Wang24,BMA26}. We therefore make no novelty claim for midpoint relations, tropical positivity, or graph encodings in isolation. The issue here is their necessity inside one fixed real measurement kernel and their exact compatibility with every stage of facial exposure, together with converse constructions at unchanged rank. The occupied off-face obstruction concerns this simultaneous compatibility, which scalar valuation inequalities alone do not test.

Positive reductions preserving singularity degree have precedent in symmetry reduction \cite{HSW23}. Our positive-retraction argument is used for a specific purpose: it transfers a real realizing kernel to every simple cone of the same rank without allowing additional ambient stresses to shorten its greedy chain. The standard Jordan-algebraic ingredients are taken from \cite{FK94}; all asymptotic and realization arguments needed for the conclusions are proved here.

\subsection{Scope and organization}
All measurements are real-linear, including those on complex and quaternionic Hermitian spaces. Error estimates control all ambient cone competitors, not only the pencils used to demonstrate sharpness. Fine central-spectrum statements concern singleton sections. A gauge reduction extends the coarse exponent and facial-profile conclusions to general nonempty affine sections on bounded test regions, but does not identify their original ungauged central curves.

We do not classify all unrestricted real facial profiles, all reducible products, arbitrary objective-barrier central paths, or interior-point iteration rates. None is needed for the three main results. The analytic proofs do not rely on finite enumeration. An optional minimum-nullity assertion, isolated in Appendix~\ref{app:computation}, is explicitly computer-assisted.

Section~\ref{sec:statements} states the three theorems. Section~\ref{sec:geometry} supplies their shared analytic foundation. Section~\ref{sec:real} proves the real one-switch construction and its Jordan and multiplicity transfers. Section~\ref{sec:fine} establishes the complete Hermitian survivor-table theorem. Section~\ref{sec:separation} proves the field separation after erasing central orders. Appendices contain minimum-rank arithmetic, reproducibility details, and precise boundary and disclosure statements.


\subsection*{AI Declaration}
Generative AI contributed to mathematical exploration, proposed constructions and proofs, exact-code development, literature searching, internal auditing, and writing. This consolidation and audit were produced with ChatGPT (GPT-6 Astra Pro). The complete model history of earlier project sessions remains to be checked by the submitting authors against the retained research record. No independent human or external referee verification is claimed. The responsible human authors must verify the mathematics and finalize this declaration before submission. A separate description of mathematical software and exact computations appears in Appendix~\ref{app:computation}.

\section{Invariants and main results}\label{sec:statements}
Let $E$ be a finite-dimensional Euclidean Jordan algebra, $e$ its unit, $N$ its rank, and $C$ its closed cone of squares. We use the trace inner product and its norm. Write $x\succeq0$ for $x\in C$. In the matrix cases $E=\Herm_n(\F)$, $\F\in\{\R,\C,\HH\}$, the inner product is $\operatorname{Re}\tr(XY)$ and $N=n$. For an idempotent $f$, $E_f$ denotes its face span and $Q_f$ its orthogonal face compression. In matrices, $Q_fX=fXf$.

Suppose $Ax=Ax_*$ has the unique solution $x_*\in C$, of rank $0<r<N$, and set $K=\ker A$. Define
\begin{equation}\label{eq:modulus}
\omega(\delta)=\sup\{\norm{x-x_*}:x\in C,\ \norm{A(x-x_*)}\le\delta\}.
\end{equation}
The sharp exponent $\alpha$ is characterized by $\omega(\delta)=\Theta(\delta^\alpha)$ as $\delta\downarrow0$. Our proof establishes existence and attainment of this exponent. Equivalent fixed norms and fixed invertible changes of measurement coordinates do not affect it.

A valid exposing stress on the current face $E_f$ is $Z\in K^\perp$ with $\ip Z{x_*}=0$ and $Q_fZ\succeq0$. The next face is the zero face of this compressed positive element inside $E_f$. A maximum-rank stress exists: a finite positive sum of available stresses has the union of their supports. Its kernel is unique. Iteration defines the \emph{greedy facial flag} and its rank drops $(s_1,\ldots,s_d)$. Lemma~\ref{lem:greedy} proves that $d$ is the minimum facial-reduction length. Put $\nu=N-r$.

The canonical center and its dual are
\begin{equation}\label{eq:center}
 c_t=\arg\max\{\log\det x:x\in\operatorname{int}C,\ x-x_*-te\in K\},
 \qquad z_t=t c_t^{-1}.
\end{equation}
Its \emph{central orders and multiplicities} are the $w_i,m_i$ for which a fixed orthogonal idempotent decomposition satisfies
\begin{equation}\label{eq:chart-statement}
 c\sum_{i=0}^k t^{w_i}c_i\preceq c_t\preceq C\sum_{i=0}^k t^{w_i}c_i,
 \quad w_0=0,\quad c_0=\supp x_*,\quad \rank c_i=m_i.
\end{equation}
The positive orders are strictly increasing, and $m_0=r$. In particular, there are $m_i$ eigenvalues of order $t^{w_i}$ and $m_i$ corresponding dual eigenvalues of order $t^{1-w_i}$. This is an order comparison, not simultaneous diagonalization or a complete Puiseux expansion.

\subsection{Harmonic certificates}
For $k\ge1$ choose parents
\begin{equation}\label{eq:harmonic}
0\le a_i<i<b_i\le k,\qquad
0=w_0<w_1<\cdots<w_k=1,\qquad
2w_i=w_{a_i}+w_{b_i}\quad(1\le i<k).
\end{equation}
For fixed parents the boundary system is nonsingular. A nonzero extreme value of a homogeneous solution propagates along increasing upper parents to the zero terminal boundary, a contradiction. Thus its solution is rational. There are at most $((k-1)!)^2$ parent assignments. Retain only those giving strict weights.

For $1\le j<k$, set
\begin{equation}\label{eq:peeling}
\ell(j)=\#\{1\le i<j:2w_i<w_{j+1}\},\qquad
u(j)=\min\{1\le p\le j:2w_p>w_j+w_{p-1}\}-1.
\end{equation}
Let $L(w)$ and $U(w)$ be the numbers of transitions in the paths beginning $k\to k-1$ and then repeatedly using $\ell$ and $u$, respectively, until zero. For $k=1$, both are one.

\begin{theoremA}[Rank-universal central and coarse spectra]\label{thm:A}
Fix an irreducible symmetric cone of rank $N$, $1\le r<N$, positive integers $m_1,\ldots,m_k$ summing to $N-r$, and $k\ge2$. The joint data consisting of the central orders $w_i$ with multiplicities $m_i$, sharp exponent $\alpha$, and minimum facial depth $d$ are attainable by a singleton section if and only if \eqref{eq:harmonic} holds and
\begin{equation}\label{eq:coarse}
 \alpha=w_v/2\le w_1\quad\text{for some }1\le v\le k,
 \qquad L(w)\le d\le U(w).
\end{equation}
For $k=1$, the orders are $0,1$, the depth is one, and $\alpha\in\{1,1/2\}$. Every allowed datum has a construction with rational coefficients in the usual real coordinates of the standard cone model.

Consequently, at fixed $N,r$, the attainable $(d,\alpha)$ spectrum is identical for all irreducible symmetric cones of rank $N$. It is the finite union of \eqref{eq:coarse} over $k\le N-r$ and all harmonic certificates, together with $(1,1)$.
\end{theoremA}

For a nonempty affine section $F=\{x\in C:Ax=b\}$ of maximum feasible rank $r$, define its bounded-set error exponent as the largest $\beta\in(0,1]$ such that
\begin{equation}\label{eq:set-error}
 \dist(x,F)\le C_B\bigl(\dist(x,C)+\norm{Ax-b}\bigr)^\beta
\end{equation}
on each bounded set for sufficiently small residuals. The coarse conclusion of Theorem A holds for these sections as well. The constants are not asserted uniform on unbounded sets. If $r=N$, the exponent is one and the depth zero; for the singleton zero target they are one and one.

Theorem A also gives an exact minimum rank defect: for permitted $d,\alpha<1$, minimize $k$ among certificates satisfying \eqref{eq:coarse}. The positive target rank is arbitrary. Appendix~\ref{app:minimum} makes this minimum finite with an explicit cutoff. This is a finite arithmetic characterization, not a claim of an elementary closed function or a polynomial-time algorithm.

\subsection{Relative flags and survivor tables}
For matrix cones write $F_i=U_0\oplus\cdots\oplus U_i$ for the ascending central flag and $V_t$ for the greedy survivor after $t$ steps. Their relative position is recorded by
\[
 x_i^t=\dim_{\F}(F_i\cap V_t)-\dim_{\F}(F_{i-1}\cap V_t),\qquad i\ge1.
\]
A simultaneous-flag congruence makes $x_i^t$ the number of grade-$i$ coordinates remaining after step $t$. This is a discrete relative-position invariant, not a specification of Euclidean angles.

An admissible table starts at $x_i^0=m_i$, ends at $x_i^d=0$, is nonincreasing in $t$, and has positive column rank drops
\begin{equation}\label{eq:margins}
 s_t=\sum_i(x_i^{t-1}-x_i^t)>0.
\end{equation}
Its first step satisfies
\begin{equation}\label{eq:terminal}
 x_k^1=0,\qquad x_i^1=m_i\quad(i<k).
\end{equation}
For $t\ge2$, define
\[
 b_t=\max\{i:x_i^{t-1}>0\},\qquad
 a_t=\min\{i:x_i^{t-1}<m_i\},
\]
and for $x_i^t>0$ put $p_t(i)=\max(\{0\}\cup\{j<i:x_j^t>0\})$. The local rules are
\begin{align}
 x_i^{t-1}>x_i^t&\ \Longrightarrow\ 2w_i\ge w_{a_t},\tag{L}\label{eq:L}\\
 x_i^t>0&\ \Longrightarrow\ 2w_i\le w_{p_t(i)}+w_{b_t}.\tag{U}\label{eq:U}
\end{align}
Rule (U) forces the entire greatest remaining grade to disappear at each step; hence $d\le k$.

\begin{theoremB}[Complete Hermitian flag spectrum]\label{thm:B}
For $\F=\C$ or $\HH$, $r\ge1$, and $k\ge2$, the central orders, positive multiplicities, relative flag table, and sharp exponent are simultaneously attainable in $\Herm_{r+\nu}(\F)$ if and only if \eqref{eq:harmonic}, \eqref{eq:margins}, \eqref{eq:terminal}, (L), (U), and $\alpha=w_v/2\le w_1$ hold. The greedy profile is exactly $(s_1,\ldots,s_d)$, and $d$ is the minimum facial depth. The construction uses no additional rank or nullity and has rational real and imaginary coefficients. For $k=1$, the only profile is $(\nu)$, with exponent one or one half.
\end{theoremB}
The same rules are necessary over $\R$, but their fixed-central-data converse is false. Taking the finite union over certificates and tables gives the complete coarse exponent/profile spectrum over $\C$ and $\HH$. It does not give the unrestricted real spectrum.

\begin{theoremC}[Field dependence after forgetting central orders]\label{thm:C}
Let $m,r\ge1$ and $\F\in\{\R,\C,\HH\}$. There is a singleton section in $\Herm_{r+m+5}(\F)$ with feasible rank $r$, sharp exponent $5/27$, and greedy facial profile $(m,1,2,2)$ if and only if
\begin{equation}\label{eq:capacity}
 r\dim_{\R}\F\ge2.
\end{equation}
The affirmative cases have rational coefficients. The negative case excludes all dense real kernels and all central certificates, not merely one prescribed spectrum. The same iff holds for general nonempty affine sections with maximum feasible rank $r$, using \eqref{eq:set-error}.
\end{theoremC}

For $m=1$, Theorem C separates real and complex rank-one problems of order seven. Real-to-complex positive embedding preserves exponent and greedy profile, so this is a strict inclusion, not just a comparison of unrelated examples. The minimum separating nullity is six under the supplementary finite verification in Appendix~\ref{app:computation}; that minimality assertion is not used in Theorem C.

\section{The shared analytic geometry}\label{sec:geometry}
We use the spectral theorem, Peirce decomposition, self-duality, and quadratic representations of Euclidean Jordan algebras \cite{FK94}. Write $Q_a=2L(a)^2-L(a^2)$, where $L(a)x=a\circ x$. If $a$ is positive and invertible, $Q_a$ is a self-adjoint cone automorphism with inverse $Q_{a^{-1}}$, and
\[
 \det(Q_ax)=(\det a)^2\det x,\qquad
 \nabla^2(-\log\det)(x)=Q_{x^{-1}}.
\]
A positive element has norm at most its trace. For mutually orthogonal idempotents $c_i$ summing to $e$, let $E_{ii}=E_{c_i}$ and $E_{ij}$ be the cross-Peirce spaces. If $a=\sum a_i c_i$, $Q_a$ acts by $a_i^2$ on $E_{ii}$ and by $a_i a_j$ on $E_{ij}$.

\begin{lemma}[Positive block bound]\label{lem:block}
If $x\succeq0$ and $s_i=\ip{c_i}x$, then $\norm{x_{ij}}\le(k+1)\sqrt{s_i s_j}$. In particular, a zero positive diagonal block forces every touching cross block to vanish.
\end{lemma}
\begin{proof}
For all $s_i>0$, apply $Q_a$ with $a=\sum s_i^{-1/2}c_i$. The result is positive and has trace $k+1$, so its norm, and hence the norm of each Peirce component, is at most $k+1$. Approximation by $x+\epsilon e$ treats zero $s_i$.
\end{proof}

\begin{lemma}[Greedy reduction has minimum length]\label{lem:greedy}
The maximum exposed support is unique at each face, and the resulting greedy chain has minimum length among all valid facial-reduction chains.
\end{lemma}
\begin{proof}
A finite positive sum of stresses realizes the union of their supports: the kernel of a sum of positive elements is the intersection of their kernels. Let $G(D)$ be the common zero face of all valid nonnegative stress functionals on a current face $D$. If $D\subset D'$, a valid stress on $D'$ restricts to a valid stress on $D$, so $G(D)\subset G(D')$. Also $G(D')$ is contained in the result of any single valid reduction from $D'$. Inductively, the greedy face after $t$ steps lies inside the face reached by any other chain after $t$ steps. The minimal feasible face belongs to every such face. A shorter chain reaching it would therefore force the greedy chain to have reached it too.
\end{proof}

\subsection{From affine sections to a singleton}
\begin{lemma}[Exact feasible-direction gauge]\label{lem:gauge}
Let $F=\{x\in C:Ax=b\}$ be nonempty, $x_0\in\ri F$, and $c=\supp x_0$. Set $M=\ker A\cap E_c$ and $\widehat A h=(Ah,P_Mh)$. Then $x_0$ is uniquely feasible for $\widehat A$. The two systems have the same sharp bounded-set exponent and exactly the same valid exposing stresses at every face containing $c$. In particular, their greedy profiles agree.
\end{lemma}
\begin{proof}
For $y\in F$, relative interior gives $x_0+\epsilon(x_0-y)\in F$ for some $\epsilon>0$. Pairing with $e-c$ and using positivity gives $\ip{e-c}y=0$, so $y\in E_c$. Conversely $x_0$ is positive definite in $E_c$, and $x_0\pm s h\in F$ for small $s$ and $h\in M$. Thus $\operatorname{aff}F=x_0+M$, and the added gauge makes the target unique.

On a bounded test region write $x-x_0=n+h$, where $n\in M$, $h\perp M$, and $\norm n\le M_0$. Put $\eta=\norm h$ and choose $\mu>0$ with $x_0\succeq\mu c$. For $x\in C$, its affine projection $z=x_0+n$ satisfies $z\succeq-\eta c$. Therefore $(1-\theta)z+\theta x_0\in F$ for $\theta=\eta/(\mu+\eta)$, and
\begin{equation}\label{eq:gauge-distance}
 \eta\le\dist(x,F)\le(1+M_0/\mu)\eta.
\end{equation}
Choose a fixed $\lambda>0$ with $(1-\lambda)\mu\ge\lambda M_0$. Then
\[
 \widehat x=x_0+\lambda h
 =\lambda x+(1-\lambda)x_0-\lambda n\in C.
\]
Its gauge residual is zero and its original residual is $\lambda\norm{Ax-b}$. Conversely, an augmented residual-$\delta$ competitor has $\norm n\le\delta$, and hence $\norm{x-x_0}\le\dist(x,F)+\delta$. The singleton residual sublevels are bounded and converge to $x_0$, as proved below; these comparisons give equality of the sharp exponents locally and on every bounded original region. Projection of an arbitrary $x$ onto $C$ changes the affine residual by at most $\norm A\dist(x,C)$, giving \eqref{eq:set-error} with the same exponent.

An augmented stress is $Z=Y+T$, with $Y\in K^\perp$, $T\in M$. On a current face containing $c$, positivity and annihilation of $x_0$ force the whole $c$-compression to vanish. Thus $0=\ip ZT=\norm T^2$, since $\ip YT=0$. Therefore $T=0$, proving the exact stress assertion. No assertion about equality of the ungauged central curves is required.
\end{proof}

\subsection{The center controls every competitor}
\begin{lemma}[Maximum-determinant envelope]\label{lem:envelope}
For a singleton section, \eqref{eq:center} exists uniquely, is semialgebraic in $t>0$, and converges to $x_*$. There is a constant $\kappa$ such that every $y\in C$ with residual at most $\delta$ satisfies
\begin{equation}\label{eq:envelope}
 0\preceq y\preceq N(1+\kappa\delta/t)c_t,
 \qquad c_t\succeq(t/N)e.
\end{equation}
\end{lemma}
\begin{proof}
Uniqueness implies $K\cap C=\{0\}$. Strict separation of $K$ from the compact trace-one section of $C$ gives $W\in\operatorname{int}C\cap K^\perp$. Pairing with $W$ bounds the positive part of every fixed shifted affine slice and every bounded-residual sublevel. The slice contains $x_*+te\succ0$. Strict concavity of $\log\det$ and its divergence to $-\infty$ at the boundary give a unique maximum. Stationarity yields $c_t^{-1}\in K^\perp$ and
\begin{equation}\label{eq:trace}
 N=\ip{c_t^{-1}}{x_*}+t\tr(c_t^{-1}).
\end{equation}
Both summands are nonnegative. Hence $c_t\succeq(t/N)e$, while $z_t=tc_t^{-1}$ has trace and norm at most $N$ and $\ip{z_t}{x_*}\le Nt$. A bounded right inverse of $A^*$ on $K^\perp$ gives
\[
 \ip{z_t}y\le Nt+\kappa N\delta.
\]
Apply $Q_{c_t^{-1/2}}$ and bound a positive element by its trace times $e$; undoing the normalization gives \eqref{eq:envelope}.

The inverse is a rational Jordan-algebraic function. Clearing its determinant in the central equations gives polynomial equations and inequalities with a unique positive solution. Their graph is semialgebraic. Thus its entries have convergent Puiseux expansions near zero; this standard semialgebraic fact is used over the fixed real coefficients of the instance \cite{BCR98}. Boundedness and uniqueness force every cluster point to equal $x_*$. The same compactness argument proves convergence of all small-residual singleton competitors.
\end{proof}

\begin{lemma}[Positive coefficient rescaling]\label{lem:rescale}
Let a finite generalized polynomial $P(t)=\sum_qt^q A_q$ have, on each block $E_{ii}$, a positive-definite first diagonal coefficient at a distinct exponent $w_i$. Suppose $(A_q)_{ij}=0$ whenever $2q<w_i+w_j$. For sufficiently large fixed $\sigma$, the polynomial $\sum_qe^{\sigma q^2}t^qA_q$ is positive definite for small $t$, and is bounded below by a positive multiple of $\sum_i t^{w_i}c_i$.
\end{lemma}
\begin{proof}
Normalize with $Q_b$, $b=\sum_i e^{-\sigma w_i^2/2}t^{-w_i/2}c_i$. The diagonal limit is positive definite on every block. A cross term survives only at $q=(w_i+w_j)/2$, when its constant multiplier is $e^{-\sigma(w_i-w_j)^2/4}$. All higher terms vanish. Choose $\sigma$ so large that the finite cross perturbation of the positive block-diagonal limit remains in the cone interior, and then choose $t$ small. The normalization gives the quantitative lower bound.
\end{proof}

\begin{proposition}[Fixed central chart and harmonicity]\label{prop:chart}
Every positive-rank singleton section has the fixed chart \eqref{eq:chart-statement}, with $0<w_1<\cdots<w_k=1$ and $k\le N-r$. Each internal grade has a nonzero cross block of the actual coefficient $H_i=A_{w_i}\in K$ joining $a<i<b$ with $2w_i=w_a+w_b$. In particular, the grades satisfy \eqref{eq:harmonic}.
\end{proposition}
\begin{proof}
Start with $c_0=\supp x_*$. Compress $c_t$ to the complementary face, take the support of its first nonzero positive Puiseux coefficient, and repeat in the remaining complementary face. This yields fixed orthogonal idempotents $c_i$ of positive ranks, strictly increasing diagonal orders $w_i>0$, and positive-definite first coefficients on $E_{ii}$. The lower eigenvalue bound in \eqref{eq:envelope} gives $w_i\le1$. Lemma~\ref{lem:block} gives the coefficient support rule
\begin{equation}\label{eq:support}
 (A_q)_{ij}=0\quad\text{if }w_i+w_j>2q.
\end{equation}
If $w_k<1$, truncate through $w_k$ and apply Lemma~\ref{lem:rescale}. All its nonconstant coefficients belong to $K$, because $c_t-x_*-te\in K$. This would yield a positive-definite exact feasible curve, contradicting uniqueness. Thus $w_k=1$.

Apply the same lemma to the truncation through one. Its residual is $O(t)$, since $A_q\in K$ for $q<1$ and $A_1-e\in K$. The envelope transfers its thick lower chart to $c_t$. The upper chart follows from the bounded trace of its positive normalization. The normalized center is semialgebraic, bounded and bounded away from the cone boundary, so it has a positive-definite limit.

For $i<k$, $H_i$ is in $K$, is positive definite on $E_{ii}$, and satisfies \eqref{eq:support}. If it had no tight cross-grade block of midpoint $w_i$, every nonzero cross-grade midpoint would be strictly smaller. Lower the whole grade $w_i$ slightly to $w'_i\in(w_{i-1},w_i)$, preserving the finitely many support inequalities, and add $t^{w'_i}H_i$ to the finite central expansion. The old support inequalities only become weaker. Lemma~\ref{lem:rescale} makes the new curve positive with residual $O(t)$ and $i$-th diagonal bounded below by $ct^{w'_i}$. This contradicts \eqref{eq:envelope} and the original $O(t^{w_i})$ diagonal. Tight cross-grade midpoints must straddle $i$. Choosing one at each grade gives \eqref{eq:harmonic}.
\end{proof}

\begin{proposition}[Sharp modulus from the restricted Hessian]\label{prop:dikin}
For $K\ne0$ put
\[
 \rho_t=\max\{\norm h:h\in K,\ \norm{Q_{c_t^{-1/2}}h}\le1\},
 \qquad
 \tau=\min_{0\ne h\in K}\max_{h_{ij}\ne0}(w_i+w_j)/2.
\]
Then $\omega(\norm{Ae}t)\asymp t+\rho_t$, $\rho_t=\Theta(t^\tau)$, and $\alpha=\min(1,\tau)$. For $K=0$, $\alpha=1$. If $k\ge2$, $\alpha=w_v/2\le w_1$ for some $v$; if $k=1$, $\alpha\in\{1,1/2\}$.
\end{proposition}
\begin{proof}
An admissible $h$ gives $c_t\pm h\in C$ with the same residual as $c_t$, so at least one has distance at least $\norm h$ from $x_*$. The point $x_*+te$ gives a linear lower bound. Conversely, a residual-$O(t)$ competitor can be written $y=c_t+h+E$, with $h\in K$ and $\norm E=O(t)$, using a bounded right inverse of $A$. The envelope bounds the normalized positive $y$, and $c_t\succeq(t/N)e$ bounds the normalization of $E$. Thus $\norm{Q_{c_t^{-1/2}}h}=O(1)$. Apply the same argument to $c_t-x_*-te\in K$, obtaining the upper bound $O(t+\rho_t)$.

Write $a_t=\sum_i t^{w_i/2}c_i$ and $c_t=Q_{a_t}M_t$, where $M_t$ stays in a compact subset of the cone interior. Differentiating the determinant identity gives Hessian covariance, and hence
\begin{equation}\label{eq:hessian}
 \ip h{Q_{c_t^{-1}}h}\asymp\norm{Q_{a_t^{-1}}h}^2
 =\sum_{i\le j}t^{-(w_i+w_j)}\norm{h_{ij}}^2.
\end{equation}
This uses covariance, not an unproved order-monotonicity assertion for $Q_a$. The finite support threshold $\tau$ is attained. A fixed vector attaining it gives one norm estimate. Projection of $K$ onto entries of weight at least $\tau$ is injective, since a vector in its kernel would have smaller maximum support weight. This finite-dimensional injectivity gives the opposite estimate uniformly in $K$, proving $\rho_t=\Theta(t^\tau)$.

For $k\ge2$, the actual nonzero coefficient $H_1$ has maximum support weight $w_1$, so $\alpha\le w_1$. Equality is a half-weight by harmonicity of the lowest grade, whose lower parent is zero. If $\alpha<w_1$, a minimizing kernel vector has no null--null block. It cannot be supported entirely on $c_0$, since small signed perturbations there preserve positivity and violate uniqueness. Its largest support-to-null block gives $\alpha=w_v/2$. For $k=1$, the only entry weights outside $c_0$ are $1/2$ and one.
\end{proof}

\begin{proposition}[Primal and dual facial bounds]\label{prop:sandwich}
The first maximum exposure is exactly the terminal grade, and every singleton section satisfies $L(w)\le d\le U(w)$.
\end{proposition}
\begin{proof}
A first PSD stress annihilating $x_*$ has zero source row. Pair successively with $H_1,\ldots,H_{k-1}$. Once lower rows vanish, the only possible pairing of $H_i$ is its positive-definite $i$-th diagonal against a positive diagonal stress block; higher--higher entries are forbidden by \eqref{eq:support}. Therefore this block, and then its row, vanish. Only the terminal grade can be exposed. The constant coefficient of $z_t=t c_t^{-1}$ is positive definite there and zero elsewhere, so that exposure is attained.

More generally, if a current face contains the grade prefix through $j$, every $H_i$ with $i\le\ell(j)$ is supported in that prefix because $2w_i<w_{j+1}$. Successive pairings force the prefix through $\ell(j)$ to survive the next step. This remains valid when the whole current face is not a grade prefix. Induction gives $d\ge L(w)$.

By the fixed chart, the dual diagonal orders are $1-w_i$ with positive-definite first coefficients, and the cross order is at least $1-(w_i+w_a)/2$. Every full dual coefficient lies in $K^\perp$. On the prefix through $j$, set $p=u(j)+1$ and truncate through $E=1-w_p$. All rows below $p$ vanish: a lower--upper cross order is at least $1-(w_{p-1}+w_j)/2>E$, and lower--lower orders are larger too. All suffix diagonal leading coefficients are included. Lemma~\ref{lem:rescale}, now for the dual orders, produces a positive-definite suffix and a zero lower prefix. This is a fixed linear combination of actual ambient dual coefficients, not an unjustified zero extension. All orders used are below one, so the stress annihilates $x_*$. Repeating gives a valid chain of length $U(w)$.
\end{proof}

\subsection{A common verification lemma for the constructions}
The following compact argument will replace repeated estimates for each realizing pencil.
\begin{lemma}[Graded positive separation]\label{lem:realization}
Let $D_t$ be a fixed diagonal power matrix with weights in $[0,1]$, zero exactly on the positive target support. Let $H_g$ be a basis of a matrix kernel $K$, with assigned positive degrees $\gamma_g$, such that
\[
 t^{\gamma_g}D_t^{-1}H_gD_t^{-1}\longrightarrow B_g
\]
exist and are linearly independent. Suppose there is $W\succ0$ orthogonal to every $B_g$. Then every PSD residual-$t$ competitor satisfies $X\preceq C D_t^2$, and its coefficient on $H_g$ is $O(t^{\gamma_g})$, up to an $O(t)$ residual correction. A thick PSD curve of residual $O(t)$ supplies the converse central-chart bound. If that curve has displacement $\Theta(t^\alpha)$ and all $\gamma_g\ge\alpha\le1$, the sharp exponent is $\alpha$.
\end{lemma}
\begin{proof}
The normalized kernels converge in the Grassmannian to $\spanof\{B_g\}$, with fixed dimension. Write $X=X_*+H+E$, $\norm E\le Ct$. Then $D_t^{-1}(X_*+E)D_t^{-1}$ is bounded. If a sequence of normalized positive competitors were unbounded, divide by its norm and take a subsequential limit. The result would be a nonzero positive element in the limiting kernel, contradicting strict positivity of $W$. Hence normalized competitors are uniformly bounded. The normalized basis has smallest singular value bounded away from zero, so its coefficients, namely the old coefficients divided by $t^{\gamma_g}$, are bounded. A thick curve is bounded above by a constant multiple of the actual center by Lemma~\ref{lem:envelope}; this gives the central lower chart. The other assertions follow from the coefficient bounds and the matching displacement.
\end{proof}

\section{Real realization and rank universality}\label{sec:real}
The converse in Theorem A is proved first with one real coordinate per positive grade and a rank-one target. A multiplicity lift and a positive Jordan retraction then give the full statement.

\subsection{One switch fills the depth interval}
Both $\ell(j)$ and $u(j)$ are nondecreasing and $0\le\ell(j)\le u(j)<j$. Indeed, if $2w_i<w_{j+1}$, tight parents give $b_i\le j$ and
$2w_i\le w_{i-1}+w_j$; hence $i\le u(j)$. The other assertions follow directly from the thresholds in \eqref{eq:peeling}.

Set $a(0)=0$ and $a(j)=1+a(\ell(j))$. Induction gives that $a$ is nondecreasing. Write the entire upper path as $j_0=k$, $j_1=k-1$, $j_{t+1}=u(j_t)$, $j_D=0$, where $D=U(w)$. For $1\le s\le D-1$, follow its first $s$ transitions and then use the lower recurrence. The resulting length is
\[
 d_s=s+a(j_s),\qquad d_1=L(w),\qquad d_{D-1}=U(w).
\]
Since $a(j)-1=a(\ell(j))\le a(u(j))\le a(j)$,
\begin{equation}\label{eq:one-switch}
 d_{s+1}-d_s\in\{0,1\}.
\end{equation}
Every integer in the interval is therefore a one-switch length. This does not assume that adding one edge changes depth by at most one.

Fix such an $s$, and put $q=j_s>0$, $J=j_{s-1}>q$. Call $i>q$ early and $i\le q$ delayed. We work on coordinates $0,\ldots,k$ with $X_*=E_{00}$ and $F_{ab}=e_ae_b^T+e_be_a^T$. Fix tight parents $P_i=(a_i,b_i)$ and a sufficiently small positive rational $\lambda$; $1/[100(k+1)^2]$ suffices. Tight physical edges are distinct, since their midpoint weights determine their children.

\subsection{Release edges, blockers, and ownership}
For each lower transition $j\to\ell(j)$, assign to its children $\ell(j)<i\le j$ the release
\[
 R_i=(0,j+1).
\]
It has midpoint at most $w_i$. For $i<j$ this is the definition of $\ell(j)$, and for $i=j$ it follows from a tight upper parent of $j$. Early children have no release.

For early $i>q$, define a minimal safe blocker
\[
 B_i=(i-1,c_i),\qquad c_i=\min\{b>i:w_{i-1}+w_b\ge2w_i\}.
\]
For delayed $2\le i\le q$, define instead
\begin{equation}\label{eq:common-blocker}
 B_i=(i-1,J).
\end{equation}
There is no additional blocker for child one. All blockers are safe: the minimal one exists by harmonicity; for \eqref{eq:common-blocker}, either $J=k$ or $q=u(J)$, both of which give $w_{i-1}+w_J\ge2w_i$. Their lower endpoints are distinct, so the blockers are pairwise distinct and non-source.

Start with $E_{ii}-E_{00}+\lambda F_{P_i}$ and add $\lambda F_{R_i}$ when a distinct release exists. If $B_i=P_i$, no separate blocker is needed. If $B_i=P_j$ for $j\ne i$, add $E_{ii}-E_{00}$ to the \emph{existing} tight generator at $j$. Safety implies $j>i$. Otherwise add $N_i=E_{ii}-E_{00}+\lambda F_{B_i}$. Each tight generator has at most one extra child, denoted $c(i)$ when present. Thus
\begin{equation}\label{eq:real-generators}
 H_i=\sum_{a\in S_i}E_{aa}-|S_i|E_{00}
       +\lambda F_{P_i}+\lambda F_{R_i},
 \qquad S_i=\{i\}\text{ or }\{i,c(i)\},
\end{equation}
with absent and duplicate releases omitted. Each main edge has one owner, not two independent coefficients.

If $\alpha=w_v/2<w_1$, also add $G=\lambda F_{0v}$. This cannot be a tight edge or a blocker. If it equals a release, the corresponding step has $j=v-1$ and $\ell(j)=0$, since $w_v<2w_1$. Omit the release from \emph{all} tight generators in that final layer. A replacement final stress is constructed below. Let $K$ be the span of the displayed generators and measure $K^\perp$. Their traces and terminal diagonals vanish, so $I,E_{kk}\in K^\perp$.

\begin{lemma}[Collision synchronization]\label{lem:collisions}
The construction has the following properties.
\begin{enumerate}[label=(\roman*),leftmargin=2em]
\item If an early blocker $B_i=P_j$ with $j>i$, its children $i,j$ belong to the same upper layer.
\item If a delayed blocker $B_i=P_j$ with $j>i$, then $\ell(q)<j\le q$ and the edge has upper endpoint $J>q$.
\item If a tight edge $P_i=(0,b)$ equals a lower release from $j=b-1$, then $i$ lies in that very lower layer and its own release equals $P_i$. A source-tight generator has no extra child; any other generator sharing that release has strictly larger weight.
\item Child one's tight source edge and every delayed source release are active throughout all noninitial upper-prefix steps.
\end{enumerate}
\end{lemma}
\begin{proof}
For (i), suppose $j$ is removed from prefix $T$ while $i$ survives. Then $i<p(T):=u(T)+1$, so $c_i\le T$. Thus
\[
2w_j=w_{i-1}+w_{c_i}\le w_{p(T)-1}+w_T<2w_{p(T)}\le2w_j,
\]
a contradiction. Since $j>i$, a suffix schedule cannot remove $j$ later.
For (ii), $j\le q$ is automatic when $s=1$. Otherwise $q=u(J)$, and $j>q$ would give $2w_j\le w_q+w_J<2w_{q+1}\le2w_j$. Also $2w_j=w_{i-1}+w_J\ge w_{q+1}$, giving $j>\ell(q)$.
For (iii), $2w_i=w_b$ and $i<b$ give $\ell(b-1)<i\le b-1$. Non-source blockers cannot add a child to a source-tight generator. A distinct sharing child has strictly larger weight because its release midpoint is at most its weight.
For (iv), when $s\ge2$, $q=u(J)\ge1$ gives $2w_1\le w_J$, so the upper endpoint of $P_1$ is at most $J$. Every delayed release has upper endpoint at most $q+1\le J$. Earlier prefixes contain these endpoints. For $s=1$ there is no such step.
\end{proof}

\subsection{Sharp recovery and the entire central chart}
Assign degree $w_i$ to $H_i$, degree $\eta_i=(w_{i-1}+w_{b(B_i)})/2\ge w_i$ to $N_i$, and degree $\alpha$ to $G$. The normalized leading form of $H_i$ is $E_{ii}+\lambda F_{P_i}$, with any additional tight release retained; an extra child diagonal vanishes because it has smaller weight. The form of $N_i$ is either $E_{ii}+\lambda F_{B_i}$ or just $\lambda F_{B_i}$, according as $\eta_i=w_i$ or $\eta_i>w_i$. Each tight pivot is unique in the leading forms by Lemma~\ref{lem:collisions}(iii), and all other main pivots are unique. Thus these limits are independent.

Set $z_i=e^{\sigma w_i^2}$, $W_{ii}=z_i$, and $W_{P_i}=-z_i/(2\lambda)$. At a new blocker of midpoint $w_i$ use the same assignment; at a strictly safe blocker use zero. All other entries, including additional releases and $G$, are zero. There is no conflicting assignment: an additional release of the same degree could equal another tight edge only if its owner had the same strict weight, in which case it was already counted once. Every leading form pairs to zero with $W$. Each nonzero normalized offdiagonal is a fixed multiple of
\[
 \exp\{-\sigma(w_a-w_b)^2/4\},
\]
so $W\succ0$ for large fixed $\sigma$. Lemma~\ref{lem:realization} gives the uniform full-ambient upper bound.

For sharpness activate only tight generators:
\begin{equation}\label{eq:real-arc}
 X(t)=E_{00}+\sum_{i=1}^{k-1}t^{w_i}H_i+tE_{kk}
       +\begin{cases}t^\alpha G,&\alpha<w_1,\\0,&\alpha=w_1.\end{cases}
\end{equation}
Every internal diagonal is $t^{w_i}$ plus nonnegative higher-order terms, and the source tends to one. After diagonal normalization, limiting offdiagonals are small multiples of $\lambda$. Strict diagonal dominance proves $X(t)\succ0$ and $X(t)\succeq cD_t^2$. Its residual is exactly $t$ and its displacement is $\Theta(t^\alpha)$. Thus the target is unique, the exponent is sharp, and the actual central chart has precisely the prescribed grades. This proof does not infer positivity of the full matrix from its two-by-two minors.

\subsection{Ambient stresses and exact minimum depth}
At the full face, source annihilation gives a zero source row. If an internal stress diagonal were positive, choose its smallest index. Its tight lower parent, source release, and any smaller extra child all have zero rows or diagonals, so pairing with its own generator is impossible. The first maximum exposure is therefore exactly $e_k$, attained by $E_{kk}$.

At a later upper prefix $0,\ldots,T$, put $p=u(T)+1$. No index below $p$ can have a positive stress diagonal. Choose the smallest one if it exists. Child one is forbidden by Lemma~\ref{lem:collisions}(iv). For $i\ge2$ its blocker is active: its upper endpoint is at most $T$, and its lower endpoint $i-1$ has zero row. A separate blocker equation forces $Z_{ii}=0$. If the blocker belongs to a tight owner, both children are still on the face by (i)--(ii), its source release is active, and a positive sum of nonnegative child diagonals would have to vanish. This is again impossible.

To attain the whole suffix $C=\{p,\ldots,T\}$, put positive diagonals $z_i=e^{\sigma w_i^2}$ on $C$ and zero elsewhere, and assign
\begin{equation}\label{eq:upper-stress}
 Z_{P_j}=-\frac{1}{2\lambda}\sum_{i\in S_j\cap C}z_i,
 \qquad Z_{B_i}=-\frac{z_i}{2\lambda}\one_{i\in C}
\end{equation}
for new blockers. Set releases and $G$ to zero. A release has upper endpoint at most $q+1$, so cannot coincide with an edge whose early owner is exposed; a delayed extra child cannot be carried by an early owner. Hence every generator pairing is zero. If an assigned edge is active and has a child in $C$, both endpoints belong to $C$: otherwise safety contradicts $w_{p-1}+w_T<2w_p$. For every contributing child,
\begin{equation}\label{eq:convex-square}
 (w_a^2+w_b^2)/2-w_i^2\ge(w_b-w_a)^2/4>0.
\end{equation}
Large $\sigma$ makes the compression positive definite on precisely $C$. Thus the first $s$ greedy steps are exactly the upper prefix and reach $0,\ldots,q$.

For an ordinary lower step $j\to\ell(j)$, put diagonal one on $C=\{\ell(j)+1,\ldots,j\}$ and zero elsewhere. Give the current release $f=(0,j+1)$ entry $-1/(2\lambda)$ and every other release zero. On a tight edge not equal to a retained release put
\[
 Z_{P_i}=-\one_{c(i)\in C}/(2\lambda),
\]
with zero numerator when no extra child exists. Give a new blocker's edge $-1/(2\lambda)$ exactly when its child is in $C$. The leading edge is zero. Own diagonals are canceled by the current release, or the identical source-tight anchor; extra diagonals are canceled by their owner edges. Lemma~\ref{lem:collisions}(iii) makes all shared assignments consistent. All nonzero offdiagonals are off-face: current releases touch $j+1$, delayed blockers touch $J>q$, and delayed owner collisions have that same off-face endpoint. Early owner children were removed together and contribute zero. The current compression is exactly the diagonal identity on $C$.

If the optional $G$ collided with the final release, that release was omitted. On the last prefix $j=v-1$, use positive diagonals $z_i=e^{\sigma w_i^2}$ on all remaining null coordinates, put
\[
 Z_{P_i}=-(z_i+z_{c(i)})/(2\lambda)
\]
with zero values outside the current block, and put $-z_i/(2\lambda)$ on current new blockers. Other releases and $G$ are zero. Earlier release owners have zero own diagonal and no extra child. Delayed extra-child corrections remain off-face. No nonzero active tight anchor touches the source: such an anchor would give $2w_i\le w_j<w_v<2w_1\le2w_i$. All pairings vanish, and \eqref{eq:convex-square} makes the final compression positive definite. This handles every leading half-weight at rank one.

Finally, every tight $H_i$ is supported on entries of weight at most $w_i$ and has positive $i$-th diagonal. If $i\le\ell(j)$, it is supported in the prefix through $j$. The same successive-pairing argument as Proposition~\ref{prop:sandwich} forces that prefix through $\ell(j)$ to survive one more step, even for rotated intermediate faces. From the greedy state $q$, at least $a(q)$ more steps are necessary. The explicit lower chain attains $a(q)$. By Lemma~\ref{lem:greedy}, the total minimum depth is exactly $s+a(q)$, completing the scalar real converse.

\subsection{Multiplicities and other irreducible cones}
\begin{lemma}[Multiplicity lift]\label{lem:multiplicity}
The scalar construction can be replaced by one with arbitrary positive grade multiplicities $m_i$, preserving the sharp exponent, complete central orders, and grouped greedy schedule. Its rank drop on a layer $C_t$ becomes $\sum_{i\in C_t}m_i$.
\end{lemma}
\begin{proof}
Replace $E_{ii}$ in every null diagonal by $I_{m_i}$, retain every physical edge on chosen group representatives, and subtract the new total diagonal trace from $E_{00}$. Observe every omitted coordinate. Choose $\lambda$ small relative to the new order. The independent main pivots are unchanged. In the leading separator use $z_iI_{m_i}$ and tight-edge entry $-m_i z_i/(2\lambda)$; the fixed multiplicity factors are dominated by the same convex-square exponential. Lemma~\ref{lem:realization} and the lifted version of \eqref{eq:real-arc}, with $tI_{m_k}$, give the full-ambient exponent and exact thick chart.

For stress equations replace $Z_{ii}$ by the trace of its group compression. A positive block with zero trace has its entire row zero, so the smallest-positive-child and forced-survivor arguments apply to arbitrary stresses and prevent premature partial exposure. To realize a scalar stress with diagonal trace $z_i$, put $(z_i/m_i)I_{m_i}$ on its group and retain its edge entries. All pairings are the old pairings. Ordinary lower stresses remain diagonal on the face. For upper and final stresses, diagonal normalization only introduces fixed factors $\sqrt{m_a m_b}$; increasing $\sigma$ in \eqref{eq:convex-square} restores positivity on every whole intended group. This proves both maximality and attainability. Fully observed positive support padding gives any $r\ge1$.
\end{proof}

\begin{lemma}[Rank-preserving positive retraction]\label{lem:retraction}
Every simple Euclidean Jordan algebra of rank $N$ contains a unital real-symmetric subalgebra $S\cong\Sym^N$ preserving rank, trace, and determinant. Its orthogonal projection $\pi$ is positive and commutes with $Q_f$ for every idempotent $f\in S$. Embedding a real kernel in $S$ and observing its entire orthogonal complement preserves the sharp exponent, central curve, and all greedy facial ranks.
\end{lemma}
\begin{proof}
Use the classification of simple algebras \cite{FK94}. In real, complex, and quaternionic Hermitian matrices take real symmetric entries; in $\Herm_3(\mathbb O)$ take real symmetric three-by-three matrices; in a rank-two spin factor use two spatial directions; rank one is $\R$. Real primitive idempotents stay primitive, and the determinant and trace restrict correctly. For $x\in C$, $y\in S_+$, $\ip{\pi x}y=\ip xy\ge0$, so $\pi x\in S_+$ by self-duality. The self-adjoint $Q_f$ preserves $S$, hence commutes with $\pi$.

An ambient residual-$\delta$ competitor has real projection satisfying the old residual bound and perpendicular component $O(\delta)$. The exponent is preserved because $\alpha\le1$; embedded old sharp competitors give the converse. The shifted affine slice lies entirely in the unital $S$, so the central curve itself is the old one.

At an old greedy face supported by $f\in S$, let $p$ be its maximal exposed support. For a new positive compressed stress $T$, $\pi T$ is an old valid stress, hence supported on $p$. Therefore $\ip{f-p}T=\ip{f-p}{\pi T}=0$. Positivity forces $T$ to be supported on $p$, including cross-Peirce components. No larger exposure is possible; the old stress attains $p$. Induction preserves the entire greedy chain, and Lemma~\ref{lem:greedy} preserves its minimum length.
\end{proof}
Propositions~\ref{prop:chart}--\ref{prop:sandwich}, the one-switch construction, and Lemmas~\ref{lem:multiplicity}--\ref{lem:retraction} prove Theorem A. The coarse affine-section conclusion follows from Lemma~\ref{lem:gauge}. Unused null directions can instead be fully observed: they join the first step and do not alter exponent or depth. The $k=1$ boundary cases are fully observed singular targets or an invisible support-to-null entry with all null diagonals observed.

\section{The complete complex and quaternionic flag spectrum}\label{sec:fine}
We prove Theorem B. The necessity is valid over all three matrix fields; only the two-coordinate construction uses $\C$ or $\HH$.

\subsection{Simultaneous flags without changing the central orders}
\begin{lemma}[Congruence normal form]\label{lem:flags}
For a matrix singleton section, one fixed congruence makes the central and greedy facial flags coordinate flags, preserving all central orders and multiplicities, the sharp exponent, and the entire greedy rank-drop profile. The relative position is precisely the table in Section~\ref{sec:statements}.
\end{lemma}
\begin{proof}
For the ascending central flag $F_i$ and the reversed survivor flag $G_j$, work modulo their common initial source space. In each $F_i\cap G_j$ choose a complement $E_{ij}$ of
\[
 (F_{i-1}\cap G_j)+(F_i\cap G_{j-1}).
\]
Induction in $i+j$ gives $F_i\cap G_j=\bigoplus_{a\le i,b\le j}E_{ab}$. A basis of all these complements is adapted to both flags; this argument works over the division rings $\R,\C,\HH$. Lift the basis and include the source. Its matrix $S$, with columns grouped by central grade, preserves each $F_i$ and is upper block triangular. Then $T=S^{-1}$ preserves those prefixes and makes every survivor coordinate.

Let $C'_t$ be the identity-shift center for $K'=TKT^*$, $X'_* = TX_*T^*$, and $A'(Y)=A(T^{-1}YT^{-*})$. The pullback $T^{-1}C'_tT^{-*}$ has original residual $tA(T^{-1}T^{-*})=O(t)$. The original envelope bounds it by a constant multiple of $C_t$. Conversely $TC_tT^*$ has transformed residual $tA(I)=O(t)$, so the new envelope bounds it by a constant multiple of $C'_t$. Thus
\begin{equation}\label{eq:congr-center}
 C'_t\asymp TC_tT^*.
\end{equation}
The identity shift is not assumed to commute with congruence. Upper triangularity gives
\[
 (D_t^{-1}TD_t)_{ab}=t^{(w_b-w_a)/2}T_{ab}\quad(a\le b),
\]
with invertible constant diagonal blocks. This matrix and its inverse are uniformly bounded. Hence $TD_t^2T^*\asymp D_t^2$, and \eqref{eq:congr-center} preserves the central chart and its flag.

Residuals correspond exactly under the congruence and norms are equivalent, preserving $\alpha$. Stresses transform by $Z'=T^{-*}ZT^{-1}$. Their restrictions to $TV$ are pullbacks of the old quadratic forms on $V$, so positivity, kernel, and compressed rank are preserved on every current face. This identifies the greedy flags and all ranks. A final source-only congruence normalizes the target to $\operatorname{diag}(I_r,0)$ without affecting the coordinate survivors.
\end{proof}

\subsection{Necessity of the two table inequalities}
The central chart and half-weight rule follow from Section~\ref{sec:geometry}. The first exposure is the entire terminal grade. At a later coordinate face let $\mu$ be the smallest weight already removed. If $2w_i<\mu$, the actual coefficient $H_i=A_{w_i}\in K$ has no entry touching a removed coordinate. All its own and lower-grade coordinates are still present. Successive pairing with the positive leading coefficients forces their rows to vanish in every PSD exposing stress. No part of grade $i$ can be removed, proving (L).

For (U), let $V$ be the current face, $W$ its greedy survivor, and $C=V\cap W^\perp$ the actual exposed subspace. Suppose a surviving grade $i$ has $2w_i>M+w_p$, where $M$ is the greatest current weight and $w_p$ the greatest strictly smaller surviving weight (zero when none). Compress the actual dual series to $V$ and truncate through $1-w_i$. Every lower surviving row is zero throughout the truncation, since its cross order with any current coordinate is at least
\[
 1-(w_p+M)/2>1-w_i.
\]
The first positive-definite diagonal coefficients on the surviving grades of weight at least $w_i$ are all included. Lemma~\ref{lem:rescale} gives a genuine ambient stress $T$ whose compression is positive definite on that high surviving subspace and zero on the lower surviving rows. This remains true for a partially surviving grade, since a principal submatrix of its positive-definite leading block is positive definite.

Let $S$ be a current stress positive definite on $C$. The entries of $T$ within $C$ and between $C$ and the high survivor may be arbitrary. Nevertheless $T+\lambda S$ is positive definite on their sum for sufficiently large fixed $\lambda$: the positive block on the high survivor remains fixed, while the Schur correction from the increasingly positive $C$ block tends to zero. All lower surviving rows remain zero. This exposes more than $C$, contradicting maximality. Thus (U) holds. In particular the greatest current grade cannot survive even partially.

\subsection{Release and blocker coordinates}
For the converse, label the null coordinates in increasing weight order, breaking ties by label, and realize the table as an ordered partition $C_1,\ldots,C_d$. Let $S_t$ be the survivor set after step $t$, including a selected source coordinate $0$, and let $t(a)$ be the removal time of coordinate $a$. Pick one representative in each grade. A tight physical pair $P_i$ joins the representatives of the parent grades in \eqref{eq:harmonic}; distinct grades have distinct tight pairs.

For $t\ge2$, choose a previously removed coordinate $q_t$ of minimum weight, breaking ties by earliest removal and then label. Keep the same choice while the minimum does not decrease. For $a\in C_t$ define
\begin{equation}\label{eq:selfrelease}
 R_a=\begin{cases}
 \{0,q_t\},&\theta_a<\theta_{q_t},\\
 \{q_t,a\},&\theta_a\ge\theta_{q_t}.
 \end{cases}
\end{equation}
The second case is a self-release. Every release is off-face at time $t$ and has midpoint at most $\theta_a$, by (L) in the first case and directly in the second.

Let $p_a$ be the immediate predecessor of $a$ in the total order restricted to $S_{t-1}$, allowing predecessor zero. Let $b_a$ be a fixed representative of the greatest-weight grade in $S_{t-2}$. Put
\begin{equation}\label{eq:fineblocker}
 B_a=\{p_a,b_a\}.
\end{equation}
The top-grade conclusion above implies that $b_a$ was removed at step $t-1$ and $\theta_{p_a}\le\theta_a<\theta_{b_a}$. If $p_a$ has the same weight as $a$, the blocker is strictly safe. Otherwise (U) at step $t-1$ gives $\theta_{p_a}+\theta_{b_a}\ge2\theta_a$; at $t=2$ the same assertion follows from harmonicity and the largest terminal weight. It is active at every earlier step and off-face at time $t$.

Blocker physical pairs are distinct. At one time the predecessor map is injective; at different times their greater-weight endpoints have strictly decreasing weights. A self-release pair is unique, because its earlier-removed endpoint cannot later exchange its role with the child. A source-release pair is assigned to only one stage: when it is used, that stage removes a strictly smaller weight, so the minimum removed weight subsequently decreases. It can, and may, be shared among children of that same stage.

For $u<v$ write
\[
 R_{uv}=e_ue_v^*+e_ve_u^*,\qquad
 I_{uv}=i e_ue_v^*-i e_ve_u^*.
\]
These are orthogonal real entry coordinates of squared norm two. Over $\HH$, fix one imaginary unit and observe all other components. Assign tight edges to real coordinates and releases to imaginary coordinates. Assign a blocker's real coordinate when its physical pair is not tight; otherwise use its imaginary coordinate if it is not a release pair.

\begin{lemma}[The only unavailable blocker is redundant]\label{lem:triple}
A blocker cannot coincide with a self-release pair. If it coincides with both a tight pair and a source release, its own tight and release pairs are that pair, and its separate blocker can be omitted.
\end{lemma}
\begin{proof}
The higher-weight endpoint of a blocker was removed before its child, while its lower endpoint survives to that child's stage. In a self-release the lower-or-equal-weight endpoint is removed first. These orders are incompatible; equal endpoint weights cannot form a blocker.

For a triple collision write $B_a=\{0,q\}$, with $q=b_a$ removed at step $t(a)-1$. A source release using $q$ can occur only at a stage $u\ge t(a)$. If $u>t(a)$, the smaller-weight coordinate $a$, removed at $t(a)$, would have decreased the historical minimum before $u$, a contradiction. Thus $u=t(a)$. Blocker safety and (L) give $\theta_a\le\theta_q/2\le\theta_a$. The tight owner is therefore in the same grade as $a$, all of whose coordinates use that same tight pair. Hence $P_{\mathrm{grade}(a)}=B_a=R_a$. Before time $t(a)$, $q$ is still present and both coordinates of that source pair vanish in any PSD stress annihilating the source. Its own tight-generator equation already forces the child's diagonal to vanish, so the extra blocker is redundant.
\end{proof}
Every nonomitted blocker now has a unique real entry coordinate absent from every tight-real or release-imaginary coordinate. There is no identification of independent coefficients.

\subsection{The rational pencil and its sharp exponent}
For every internal coordinate $a$, use
\begin{align}
 H_a&=E_{aa}-E_{00}+\lambda R_{P_{\mathrm{grade}(a)}}+\lambda I_{R_a},\label{eq:fineH}\\
 N_a&=E_{aa}-E_{00}+\lambda F_{B_a}\quad\text{when not omitted},\label{eq:fineN}
\end{align}
where $F_{B_a}$ is the assigned component. When $\alpha=w_v/2<w_1$, add $G=\lambda R_{0,\mathrm{rep}(v)}$. Its real coordinate is unused because its midpoint is smaller than every tight or blocker midpoint; a release there is imaginary. Take $X_*=\operatorname{diag}(I_r,0)$, $K$ the real span, and $A=P_{K^\perp}$. Choose $\lambda=1/[100(r+\nu)^2]$. All generators have trace zero and zero terminal diagonal.

Assign degrees $\theta_a$ to $H_a$, the safe midpoint to $N_a$, and $\alpha$ to $G$. Their scaled limits are independent. First eliminate the unique blocker pivots and the optional leading pivot; then each $H_a$ has its own null diagonal, even when several children share one tight edge.

Set $W_{aa}=e^{\sigma\theta_a^2}$, the source diagonals equal to one, and on a tight real pair of grade $i$ set $W_{P_i}=-e^{\sigma w_i^2}/(2\lambda)$. This is consistent for all children sharing the pair. Assign the same value at a blocker whose midpoint equals its child's weight and zero at a strictly safe blocker. All release-imaginary and leading entries are zero. Lemma~\ref{lem:triple} prevents conflicts. Every leading generator is orthogonal to $W$. A leading self-release can join equal weights, but its separator entry is zero; every \emph{nonzero} separator edge joins distinct weights and its normalized size tends to zero as $\sigma\to\infty$. Thus $W\succ0$ for large $\sigma$, also over $\HH$ by the norm bound on cross terms. Lemma~\ref{lem:realization} proves the all-competitor upper bound.

The thick sharp curve is
\[
 X(t)=X_*+\sum_{a\text{ internal}}t^{\theta_a}H_a
       +tI_{\mathrm{terminal}}+[t^\alpha G\text{ when present}].
\]
Every null diagonal has exactly its assigned power, the selected source diagonal tends to one, and every normalized edge is bounded by a fixed multiple of $(r+\nu)\lambda$. Strict diagonal dominance gives positivity and thickness. Its residual is exactly $t\norm{I_{\mathrm{terminal}}}$, since the entire terminal diagonal is measured, and its displacement is $\Theta(t^\alpha)$. The central envelope and the common verification lemma establish the exact orders, multiplicities, and sharp exponent.

\subsection{Exact maximal exposures}
Put $c=1/(2\lambda)$ and take $Z_1=I_{\mathrm{terminal}}$. For $t\ge2$ use the ambient stress
\begin{equation}\label{eq:fine-stress}
 Z_t=\sum_{a\in C_t}E_{aa}
 -c\sum_{\substack{\text{distinct release pairs }e\\\text{assigned to }C_t}} I_e
 -c\sum_{\substack{a\in C_t\\N_a\text{ present}}}F_{B_a}.
\end{equation}
Every original generator pairing is zero. A release belongs to exactly one stage, where it cancels the unit diagonal pairing of each sharing child. Each blocker coordinate is unique and cancels its child's diagonal. Tight-real and optional leading entries receive zero. All the offdiagonal entries in \eqref{eq:fine-stress} are off-face, so its compression is exactly the identity on $C_t$.

To prove maximality at an earlier stage, take the first future coordinate $a$ in the total order with positive diagonal in a purported larger PSD stress. Its blocker predecessor is another earlier future coordinate or the source, so its whole row is zero. The blocker's upper endpoint is still on-face because it is removed only at $t(a)-1$. The equation for $N_a$ forces $Z_{aa}=0$, a contradiction. If the blocker was omitted, Lemma~\ref{lem:triple} uses the still-active source-tight pair in $H_a$ for the same contradiction. This argument also applies at the initial face and kills every internal coordinate. Therefore the prescribed faces are the actual greedy faces, not just a valid chain. Lemma~\ref{lem:greedy} proves the minimum degree assertion. This completes Theorem B.

\section{Field dependence after erasing the central spectrum}\label{sec:separation}
We now prove Theorem C analytically. The issue is not whether a particular complex pencil can be made real, but whether \emph{any} real kernel with any central certificate can have the specified exponent and profile.

\subsection{Off-face saturation}
\begin{lemma}[Occupied off-face space]\label{lem:cut}
Let $f$ be a current face idempotent in a Euclidean Jordan algebra and $g\le f$ an intended exposed subface. Write $O_f=\Id-Q_f$. Suppose $A_1,\ldots,A_l,B\in K$ satisfy
\[
 Q_gA_j=0,\qquad 0\ne Q_gB\succeq0,\qquad
 O_fB\in\spanof\{O_fA_1,\ldots,O_fA_l\}.
\]
No ambient measurement stress can have positive compression supported exactly on $g$ and positive definite there.
\end{lemma}
\begin{proof}
Cancel the off-face component to obtain $J=B-\sum_jc_jA_j\in K\cap E_f$. Then $Q_gJ=Q_gB$ is nonzero positive. An alleged stress $Z\in K^\perp$, with compression $S$ positive definite exactly on $g$, would give $0=\ip ZJ=\ip S{Q_gJ}>0$.
\end{proof}

The exact linear-algebra alternative behind this obstruction is
\[
 Q_g(K\cap E_f)\cap C_g=\{0\}
 \quad\Longleftrightarrow\quad
 \exists S_g\in\operatorname{int}C_g:\ S_g\perp Q_g(K\cap E_f).
\]
Strict separation of the trace-one positive section proves the equivalence. The identity $(K\cap E_f)^\perp_{E_f}=Q_fK^\perp$ lifts its zero-extended positive separator to an actual ambient stress. Lemma~\ref{lem:cut} is the concrete positive-circuit use of this alternative.

There is a useful weighted form with an important restriction. In aligned coordinates let $K_{\le u}$ consist of kernel elements supported on entries of weight at most $u$, and put $R_f(u)=\dim_{\R}O_f(K_{\le u})$. If $\tau$ is the smallest grade in the intended exposed block and
\begin{equation}\label{eq:internal-restriction}
 0<\tau<1,
\end{equation}
then its actual central coefficient $H_\tau\in K$ has nonzero PSD compression to that block, whereas every element of $K_{<\tau}$ has zero compression there. Therefore
\begin{equation}\label{eq:cutjump}
 R_f(\tau)>R_f(\tau^-)
\end{equation}
is necessary. This is not asserted sufficient. The restriction \eqref{eq:internal-restriction} cannot be dropped: at the first terminal exposure $O_e=0$, and $A_1$ need not belong to $K$. Only $A_1-e\in K$. Appendix~\ref{app:scope} records the simplest counterexample to the unrestricted jump formulation.

\subsection{The denominator restricts every possible certificate}
Let $h=k-1$ be the number of internal harmonic grades and $Q=2I-B$ their integer midpoint matrix. Its determinant counts selections of one of the two parents at each internal vertex that contain no internal directed cycle. Indeed, determinant expansion over disjoint directed cycles equals inclusion--exclusion for the bad parent selections. Thus
\begin{equation}\label{eq:detbudget}
 1\le\det Q\le2^h,\qquad \den(2\alpha)\mid\det Q.
\end{equation}
Positivity of the determinant also follows from nonsingularity and the existence of acyclic selections. Since the first rank drop is precisely the terminal multiplicity,
\begin{equation}\label{eq:profilebudget}
 h\le\nu-s_1,\qquad \den(2\alpha)\le2^{\nu-s_1}.
\end{equation}

For $\nu=m+5$, $s_1=m$, and $\alpha=5/27$, the certificate contains $2\alpha=10/27$. Hence $h\le5$, but $27>2^4$, so $h=5$ and $\det Q=27$, the only positive multiple of 27 at most 32. Every internal grade is simple; the terminal grade has multiplicity $m$. Every internal weight is an integer divided by 27, the list contains $10/27$, and its minimum is at least $5/27$.

\begin{lemma}[Five-vertex determinant rigidity]\label{lem:27}
A strict harmonic certificate with five internal vertices, determinant 27, a weight $10/27$, and minimum weight at least $5/27$ has one of the two ordered numerator lists
\begin{equation}\label{eq:two-lists}
 A=(0,5,10,13,16,20,27),\qquad
 B=(0,5,8,10,12,16,27).
\end{equation}
\end{lemma}
\begin{proof}
Exactly five of the 32 parent selections contain a cycle. A two-cycle alone occurs in eight selections, so none exists. Distinct cycles of length at least three cannot coexist in one functional graph on five vertices. If $c_j$ counts directed $j$-cycles in the full parent graph, their selection events are disjoint and
\[
 4c_3+2c_4+c_5=5.
\]
In particular a Hamiltonian five-cycle exists. Fix one and label it $1\to2\to3\to4\to5\to1$. Each vertex has at most one additional internal edge. A chord of forward distance four would create a forbidden two-cycle. A chord of distance three creates a three-cycle, whose four bad selections, together with the Hamiltonian cycle, exhaust the budget. In this case there is exactly one such chord and every other second edge goes to a boundary. If there is no distance-three chord, every chord has distance two and creates a four-cycle costing two selections. There must be two chords; their tails must be adjacent, since nonadjacent tails create an additional three-cycle. Three or more chords already exceed the budget. Thus, up to cyclic relabeling, the two types are a single chord $1\to4$, or the two chords $1\to3$, $2\to4$.

The boundary endpoints of the remaining second edges are bits. Solving the two five-variable midpoint systems gives, in cycle order,
\begin{equation}\label{eq:bitmatrices}
27x=
\begin{pmatrix}8&4&10&5\\14&7&4&2\\1&14&8&4\\2&1&16&8\\4&2&5&16\end{pmatrix}b,
\quad b\in\{0,1\}^4;
\qquad
27x=
\begin{pmatrix}12&10&5\\9&12&6\\15&8&4\\3&16&8\\6&5&16\end{pmatrix}b,
\quad b\in\{0,1\}^3.
\end{equation}
These identities are obtained by substituting $2x_i=x_{i+1}+b_i$ at boundary rows and the displayed chord equation at the other rows; they can also be verified by direct multiplication by $Q$. In the first matrix a row has subset sum ten only when $b=(0,0,1,0)$ or $(1,0,0,1)$. The former yields a numerator four and violates the minimum barrier. The latter has sorted list $(5,10,13,16,20)$. In the second matrix the only such bit vector is $(0,1,0)$, giving $(5,8,10,12,16)$. Including the boundaries proves \eqref{eq:two-lists}. This is an analytic restriction on all parent graphs, not an enumeration of measurement kernels.
\end{proof}

\subsection{The profile forces a saturated real cut}
Apply Lemma~\ref{lem:flags}. All five internal grades are simple, and the first layer is the whole terminal grade. Rule (U) forces the next unit drop to remove the greatest internal grade, coordinate five.

For list $A$, the next layer must contain numerator 16. Pairing it with five violates (L), since $10<20$. Pairing it with ten leaves surviving numerators five and thirteen, violating (U), since $26>16+5$. Pairing it with thirteen leaves a last exposure of five, violating (L), since $10<13$. Therefore list $A$ cannot realize $(m,1,2,2)$ over any field.

For list $B$, the next layer must contain numerator 12. Pairing it with five violates (L), since $10<16$. Pairing it with eight leaves ten and five, violating (U), since $20>12+5$. The only possible schedule is consequently
\begin{equation}\label{eq:forcedschedule}
 \{\text{terminal grade}\},\quad\{5\},\quad\{3,4\},\quad\{1,2\}.
\end{equation}
At its third step the current face consists of the source and coordinates $1,2,3,4$, and the proposed exposed block is $g=\spanof(e_3,e_4)$. Let $H_2,H_3$ be the actual central coefficients of orders $8/27,10/27$. Their only possible off-face component is the source-to-5 block: the source-to-terminal weight is $1/2$, and the lightest null-to-5 weight is $21/54>10/27$.

The unique tight midpoint for $8/27$ in list $B$ is $(0,5)$, so $H_2$ has a nonzero component in that cut. Its compression to $g$ is zero. The $g$-compression of $H_3$ is a positive multiple of $E_{33}$, and its only possible off-face component is in the same cut. The cut has real dimension $r\dim_{\R}\F$. If this dimension is one, its nonzero $H_2$ component spans it, and Lemma~\ref{lem:cut} forbids the intended exposure. Equivalently, in the real rank-one case
\[
 J=H_3-\frac{(H_3)_{05}}{(H_2)_{05}}H_2
\]
is supported on the current face and has nonzero positive compression to $g$. This proves necessity of \eqref{eq:capacity}. Lemma~\ref{lem:27} shows why choosing different central orders cannot avoid it.

\subsection{A six-generator witness and its full-ambient estimate}
For $m=r=1$ over $\C$, put $\lambda=1/100$, $X_*=E_{00}$, and let $K$ be the real span of
\begin{align}
 H_1&=E_{11}-E_{00}+\lambda R_{03},&
 H_2&=E_{22}-E_{00}+\lambda R_{05},\notag\\
 H_3&=E_{33}-E_{00}+\lambda(R_{24}+I_{05}),&
 H_4&=E_{44}-E_{00}+\lambda R_{25},\label{eq:witness}\\
 H_5&=E_{55}-E_{00}+\lambda R_{16},&
 N_2&=E_{22}-E_{00}+\lambda R_{14}.\notag
\end{align}
Measure $K^\perp$. Its orthogonal projection has rational coefficients in standard real Hermitian coordinates. Each generator has trace zero and terminal diagonal zero.

Let $X\succeq0$ have residual at most $\delta$. Write $X=Y+E$, $\norm E\le\delta$, and
\[
Y=E_{00}+aH_1+bH_2+cH_3+dH_4+eH_5+fN_2.
\]
Then $Y+\delta I\succeq0$, $\tr Y=1$, and $Y_{66}=0$. Put $U_j=|Y_{jj}|+\delta$. Trace bounds give $U_0=O(1)$ and $U_6=\delta$. Principal two-by-two positivity and the generator coordinates give
\begin{equation}\label{eq:witnessineq}
\begin{split}
U_1&\le C\sqrt{U_3},\qquad
U_2\le C\bigl(\sqrt{U_5}+\sqrt{U_1U_4}\bigr),\\
U_3&\le C\sqrt{U_2U_4},\qquad
U_4\le C\sqrt{U_2U_5},\qquad
U_5\le C\sqrt{U_1\delta}.
\end{split}
\end{equation}
The additive $\delta$ terms are absorbed since $U_j\ge\delta$ and the relevant square roots dominate $\delta$ up to fixed constants. The sum in the second inequality is required: $Y_{22}=b+f$, so neither coefficient is assumed nonnegative or equal to that diagonal.

Successive substitution gives
\[
 U_1\le C\delta^{1/15}U_2^{2/5},\qquad
 U_2\le C\bigl(\delta^{4/15}U_2^{1/10}+\delta^{1/6}U_2^{1/2}\bigr).
\]
One summand controls $U_2$ up to a fixed factor. It gives respectively $U_2=O(\delta^{8/27})$ or $O(\delta^{1/3})$, both implying the former. Consequently
\begin{equation}\label{eq:witnessorders}
(U_1,\ldots,U_5)=
\bigl(O(\delta^{5/27}),O(\delta^{8/27}),O(\delta^{10/27}),
O(\delta^{12/27}),O(\delta^{16/27})\bigr).
\end{equation}
Also $|f|\le C\sqrt{U_1U_4}=O(\delta^{17/54})$ and $|b|\le C\sqrt{U_5}$. Thus all coefficients and the unstructured error are controlled, proving $\norm{X-X_*}=O(\delta^{5/27})$ for every ambient competitor.

A sharp thick curve is
\begin{equation}\label{eq:witnessarc}
X(s)=E_{00}+s^5H_1+s^8H_2+s^{10}H_3+s^{12}H_4+s^{16}H_5+s^{27}E_{66}.
\end{equation}
Its null diagonal entries are the displayed powers and its source tends to one. The normalized tight edges are small constants, while the extra imaginary entry is smaller by $s^2$. Hence $X(s)\succ0$ for small $s$. Its residual is exactly $s^{27}$ and displacement $\Theta(s^5)$. Equations~\eqref{eq:witnessorders} and \eqref{eq:witnessarc}, together with the envelope, establish the actual center's weights $B$ and sharp exponent $5/27$.

With $c=50$, ambient measurement stresses are
\begin{align}
Z_1&=E_{66},& Z_2&=E_{55}-cR_{16},\notag\\
Z_3&=E_{33}+E_{44}-c(I_{05}+R_{25}),&
Z_4&=E_{11}+E_{22}-c(R_{03}+R_{05}+R_{14}).\label{eq:witnessstress}
\end{align}
Their 24 generator pairings vanish, and the successive compressions expose $\{6\},\{5\},\{3,4\},\{1,2\}$. For maximality, a first PSD stress has zero source row; $H_1,H_2$ kill rows one and two, then $H_3,H_4$ kill rows three and four, and $H_5$ kills row five. After removing six, the same first four equations still kill rows one through four. After removing five, $H_1$ kills row one and $N_2$ kills row two. The displayed stresses attain all remaining allowed ranks. Thus the profile is exactly $(1,1,2,2)$ and its length is minimal.

The same complex-subfield witness works over $\HH$ with all other components observed. For a real rank-two witness, add positive source coordinate seven, replace $I_{05}$ by $R_{57}$ in \eqref{eq:witness} and \eqref{eq:witnessstress}, and take target $E_{00}+E_{77}$. Source annihilation and bounded positive source diagonals leave every argument unchanged. Positive support padding gives larger $r$; $m-1$ fully observed null coordinates with all added crosses observed join the first exposure and have central order one. These operations preserve the exponent and every later rank drop. All affirmative cases of Theorem C follow.

Embedding a real kernel into a complex Hermitian space and observing its imaginary complement preserves all invariants by Lemma~\ref{lem:retraction}. The separation is therefore a strict inclusion of attainable coarse profile spectra. Lemma~\ref{lem:gauge} proves the general affine-section assertion. Theorem C is complete without any computation-assisted minimality claim.

\appendix
\section{Exact minimum rank defect}\label{app:minimum}
For $\alpha<1$ and an attainable $(d,\alpha)$, Theorem A gives
\begin{equation}\label{eq:minrank}
\nu_{\min}(d,\alpha)=\min\{k:\text{a harmonic certificate has }
\alpha=w_v/2\le w_1,\ L(w)\le d\le U(w)\}.
\end{equation}
The minimum total matrix order at positive target rank $r$ is $r+\nu_{\min}$. The formula is arithmetic, since the weights solve $(2I-B)w=c$ and hence are ratios of integer determinants. Here is an explicit finite cutoff, included to distinguish \eqref{eq:minrank} from an unbounded enumeration.

If $0<\alpha<1/2$ is rational, put $e=\lceil\log_2(1/\alpha)\rceil$, $c=2^{-(e-1)}$, and write $(2\alpha-c)/(1-c)=a/b$ in lowest terms, including $0/1$. A base harmonic certificate uses the dyadic spine $1/2,1/4,\ldots,c$, the affine grid $c+j(1-c)/b$ for $1\le j<b$, and terminal one, omitting duplicates. Spine parents are zero and double the weight; new grid parents are its two grid neighbors. The vertex $2\alpha$ is present and all internal weights are at least $c\ge\alpha$. The lower peeling removes the spine levels successively, together with all intervening grid weights, so $L=e$. The real lower-endpoint construction realizes $(e,\alpha)$ with nullity at most $e+b-1$.

To realize a prescribed $d>e$, choose $M\ge1$ with
\[
(2^M-1)(1-2\alpha)>\alpha(d-2).
\]
Set $A=2^M-1$, $\Delta=(2A+d-2)^{-1}$, $\beta=A\Delta>\alpha$, $T_0=2\beta+\Delta$, and $k=M+d-1$. Define
\[
w_1=\beta,\qquad
w_i=T_0-(T_0-\beta)2^{-(i-1)}\ (2\le i\le M+1),
\qquad
w_i=2\beta+(i-M-1)\Delta\ (M+2\le i\le k).
\]
Then $w_{M+1}=2\beta$ and $w_k=1$. Parents are $(0,M+1)$ for vertex one, $(i-1,M+2)$ for $2\le i\le M+1$, and the two neighbors for the remaining internal vertices. The upper recurrence removes the top arithmetic-spine coordinates one at a time, then the $M$ middle coordinates, then vertex one. Thus $U=d$ and this helper has exponent $\beta$ at depth $d$ and nullity $M+d-1$.

Glue the base and helper only at their positive source, keeping all null coordinates and edges separate and measuring each terminal. Their generator equations decouple in every compressed stress because its source row is zero. Hence maximum exposure acts componentwise and the degree is the maximum of the two degrees. The positive leading separators combine after diagonal scaling, and a sufficiently small common edge coefficient makes the combined normalized sharp curves positive. Equivalently, Lemma~\ref{lem:realization} applies to their union, whose leading pivots outside the source are disjoint. The exponent is the smaller one, $\alpha$. Therefore the search in \eqref{eq:minrank} may be restricted to
\[
d\le k\le
B(d,\alpha):=
\begin{cases}
e+b-1,&d=e,\\
e+b+M+d-2,&d>e.
\end{cases}
\]
This cutoff is not asserted optimal; the finite minimum is exact.

For completeness, the dimension-unrestricted spectra are $\{1,1/2\}$ at depth one, $\Q\cap[1/4,1/2]$ at depth two, and $\Q\cap[2^{-d},1/2)$ at $d\ge3$. Rationality follows from the harmonic system. The guaranteed lower bound $\alpha\ge2^{-d}$ is the classical facial-reduction bound \cite{Sturm00,Lourenco21}; the constructions above attain every listed rational value. If $\alpha=1/2$, the half-weight rule requires $w_1\ge1/2$, and the upper recurrence has at most two steps. The special minima are $\nu_{\min}(1,1)=\nu_{\min}(1,1/2)=1$ and $\nu_{\min}(2,1/2)=2$. All constructions can use rational positive scales: clear the weight denominators and replace $e^{\sigma w_i^2}$ by integer powers of a sufficiently large integer. The finitely many strict convex-square inequalities used for diagonal dominance then remain strict. A substitution $t=s^Q$ clears all curve exponents.

\section{Reproducibility and the separate minimality assertion}\label{app:computation}
The proofs of Theorems A--C are analytic. The supplied standard-library exact programs test their certificates, not floating-point fits of exponents. The one-switch audit checks 225 constructions in 17 families. The general-table audit checks 7,572 labelled exponent--schedule cases, including 1,590 split-grade and 404 non-prefix cases, with 247,749 exact generator/stress pairings. These are finite checks and do not establish the all-dimensional statements.

A separate audit of Theorem C solves the determinant-27 candidate systems and verifies the explicit witness. It also compares the complete complex candidate spectrum through nullity five, as given by Theorem B, with exact real constructions. The counts of distinct $(\text{profile},\alpha)$ pairs for $\nu=1,\ldots,5$ are $2,4,11,48,292$. Each candidate is covered by a one-switch construction, a documented real tied-edge witness, or fully observed padding. Conditional on this reproducible finite enumeration, nullity six is the first field-separating nullity. This last assertion is \emph{computer-assisted}; omitting it leaves every principal theorem unchanged. The machine-readable certificates and historical witness specifications are retained with the research packet, not used as substitutes for the proofs in the text.

A newly written independent checker additionally verifies the two determinant-27 graph types on 3,125 anchored Hamiltonian graphs, solves both boundary matrices by exact elimination, and checks all 24 pairings and all face compressions of the six-generator witness. It confirms that the only permitted off-face pair at the obstructed cut is $(0,5)$. It also records the terminal exception described next. The checker is structurally independent of the supplied witness audit, but neither an automated check nor this internal audit is an external referee review.

\section{Scope boundaries, audit correction, and disclosure}\label{app:scope}
\subsection{An essential correction to the cut-rank formulation}
The weighted consequence \eqref{eq:cutjump} is only asserted for internal $0<\tau<1$. To see why, take $X_*=E_{00}$ and $K=\spanof(R_{01})$ in $\Sym^2$. The center is $\operatorname{diag}(1+t,t)$ and the first stress $E_{11}$ exposes the terminal grade one. At the initial full face, $O_f=0$, so both cut ranks are zero and no strict jump occurs. The central coefficient at order one is $I$, not an element of $K$. The unrestricted version of the weighted jump in the earlier research note was therefore false. Lemma~\ref{lem:cut} itself is valid without this restriction, and Theorem C uses it at the internal order $10/27$, so its proof is unaffected.

\subsection{What the statements do not assert}
The full unrestricted real facial-profile spectrum remains unclassified. It is not a missing hypothesis in Theorems A--C. The survivor inequalities alone are not sufficient over $\R$, and the cut-rank jumps have not been proved sufficient for all real kernels. No simple universal capacity threshold is inferred from the exact threshold for the family in Theorem C.

Irreducibility is necessary for Theorem A. An orthant admits only linear singleton recovery: its tangent cone at the target intersects the measurement kernel only at zero, and compactness of the unit tangent section gives a uniform linear residual bound. The rank-two PSD example in the introduction has exponent $1/2$. Coupled products require their Peirce support and feasible-rank distribution to be retained. Their classification and quotient-encoding questions are outside this paper.

The identity-shift central curve is not an arbitrary optimization central path. Eigenvalue orders are not full Puiseux ramification data. The gauge reduction is not claimed to preserve the ungauged central spectrum. The flag theorem is a congruence statement, not simultaneous orthogonal diagonalization or classification of Euclidean angles. Construction at rational coefficients does not imply genericity or uniform conditioning constants. These distinctions are maintained in every use of the results.

\subsection{Pre-submission status and research-process disclosure}
This consolidated text contains the complete research arguments for the stated principal results and the correction above. It is not a report of independent specialist certification or worldwide bibliographic priority. Before submission, the responsible authors must verify authorship, institutional information, the bibliography, the computational appendix, and the proofs, especially the Jordan transfer and the two edge-ownership constructions. No author names or affiliations have been supplied in this draft.

Generative AI was used in the project's mathematical exploration, proposed proofs and constructions, code development, literature searching, auditing, and consolidation of this manuscript. This draft and its internal audit were produced with ChatGPT; the exact model/tool history of earlier sessions must be verified from the project record by the submitting authors. No human proof verification or external peer review is represented as having occurred. The final submission should contain an accurate journal-compliant AI declaration rather than treating this provisional statement as a substitute for author responsibility.

\bibliographystyle{amsplain}
\bibliography{references}
\end{document}
